% optimization/part_intro-DE.tex
\chapter*{Optimierung}
\phantomsection
\addcontentsline{toc}{chapter}{Optimierung}

\section*{Kurzfassung für Leserinnen und Leser}

Optimierung hilft, wenn ein ingenieurtechnisches Problem mehr als eine
mögliche Änderung zulässt. Sie durchsucht erklärte Freiheitsgrade nach
Kandidaten, die erklärte Nebenbedingungen erfüllen und unter einem erklärten
Ziel gut abschneiden. Sie kann weder den Zweck dieses Ziels liefern noch die
Anwendbarkeit ihrer Eingaben feststellen, einen Kandidaten annehmen oder eine
Engineering Decision treffen.

Dieser Teil verfolgt ein Beispiel. Eine Vermessungspolylinie und eine
CAD-Kurve weichen in einem Übergangsbereich voneinander ab. Die Identität des
Alignments bleibt fest. Die Ingenieurin oder der Ingenieur erlaubt eine
begrenzte Änderung ausgewählter intrinsischer Krümmungsparameter; Endpunkte,
Endrichtungen, Projektkontext und Quellbeobachtungen bleiben fest. Die Frage
lautet nicht »Welche Kurve ist wahr?«, sondern:
\begin{quote}
Welche Krümmungsänderung ist für diesen erklärten Zweck, unter diesen
Nebenbedingungen, Annahmen und Beobachtungen ein prüfbarer Kandidat?
\end{quote}

\section*{Der Vertrag vor der Berechnung}

Ein Optimierungsproblem erhält erst durch erklärte Verantwortlichkeiten eine
Bedeutung. Für das durchgehende Beispiel sei
\[
\mathcal P=(I,P,M,x,q,\mathcal F,J,\Sigma),
\]
wobei \(I\) das Engineering Object identifiziert, \(P\) den Zweck nennt,
\(M\) das Rekonstruktions- und Bewertungsmodell bezeichnet, \(x\) die freien
Variablen, \(q\) feste Eingaben und Kontext, \(\mathcal F\) die zulässige
Menge, \(J\) das Ziel mit Einheiten und Skalierung und \(\Sigma\) die
Unsicherheitsgrenze enthält. Kein Eintrag darf aus einem Solver-Ergebnis
abgeleitet werden.

Der Leserpfad folgt den Abhängigkeiten dieses Vertrags:
\[
\begin{array}{c}
\text{ingenieurtechnische Frage}\\
\downarrow\\
\text{erklärte Freiheiten, feste Eingaben, Ziel, Nebenbedingungen}\\
\downarrow\\
\text{numerische Rekonstruktion und Residuenbewertung}\\
\downarrow\\
\text{mathematischer Kandidat und Konvergenzevidenz}\\
\downarrow\\
\text{Anwendbarkeitsprüfung und ingenieurtechnische Evaluation}\\
\downarrow\\
\text{Engineering Decision außerhalb des Solvers.}
\end{array}
\]

\section*{Warum die technischen Kapitel in dieser Reihenfolge folgen}

Numerische Rekonstruktion erscheint zuerst, weil jede versuchsweise
Krümmungsänderung in Geometrie überführt werden muss, bevor Endpunkt- und
Beobachtungsresiduen auswertbar sind. Die Alignment-Optimierung formuliert
danach das kontinuierliche und diskrete Problem und führt SQP, Ableitungen und
dünne Besetzung erst ein, wenn das wachsende Problem sie erfordert. Das
Architekturkapitel beschreibt die Verantwortungsgrenzen eines möglichen
Solver-Dienstes und trennt sie von der begrenzten experimentellen
AXTRAN-Evidenz. Die abschließende Spezialisierung untersucht, wie ein erklärtes
Realisierungsmodell in die Bewertung eingehen kann, ohne Modell, Realisierung,
Beobachtung oder Entscheidung zusammenfallen zu lassen.

\begin{figure}[ht]
\centering
\begin{tikzpicture}[node distance=1.9cm,box/.style={draw,rounded corners,minimum width=3.8cm,minimum height=0.9cm,align=center},arrow/.style={->,thick}]
\node[box] (params) {Parameter};
\node[box,right=of params] (eval) {Bewertung};
\node[box,below=of eval] (residual) {Residuum};
\node[box,left=of residual] (updated) {Aktualisierte Parameter};
\draw[arrow] (params)--(eval); \draw[arrow] (eval)--(residual);
\draw[arrow] (residual)--(updated); \draw[arrow] (updated)--(params);
\end{tikzpicture}

\caption{Optimierung erzeugt einen überarbeiteten Kandidaten und Evidenz für die Prüfung; die Schleife nimmt den Kandidaten nicht an.}
\label{fig:icon-optimization-evaluation-loop}
\end{figure}

\begin{principle}[Optimierungsprinzip]
\label{princ:optimization-principle}
Eine AIM-Optimierungsformulierung muss Engineering Object und Zweck,
Freiheitsgrade und feste Eingaben, Modell, Ziel und Skalierung,
Nebenbedingungen, Unsicherheitsgrenze und Provenienz erklären. Ihr numerisches
Ergebnis ist ein Kandidat mit Residuen- und Konvergenzevidenz, keine
Engineering Decision.
\end{principle}

\begin{summary}
Optimierung macht aus Evaluation keine Autorität. Sie macht eine erklärte
ingenieurtechnische Frage rechnerisch durchsuchbar und liefert Kandidaten, die
geprüft, verglichen und evaluiert werden können.

\[
\boxed{\text{berechnet}\neq\text{anwendbar}\neq
\text{angenommen}\neq\text{autoritativ}}
\]
\end{summary}
